% ============================================================
\chapter{Product Measures and Fubini-Tonelli}
\label{ch:fubini}
% ============================================================

\section{Introduction}

Computing multiple integrals is one of the most frequent operations in analysis and
probability. The theorems of Fubini and Tonelli provide conditions under which a
double integral can be computed as an iterated integral\index{iterated integral},
in either order.

\begin{intuition}[title=\textbf{Why Fubini-Tonelli?}]
Computing $\iint f(x,y)\, d(x,y)$ directly on a product is often difficult.
Fubini and Tonelli allow us to reduce this to successive integrations: first
integrate in $y$ for each fixed $x$ (or vice versa). Tonelli handles the case of
nonnegative functions (no integrability hypothesis), while Fubini applies to
integrable functions.
\end{intuition}

% ------------------------------------------------------------
\section{Product $\sigma$-algebra}
% ------------------------------------------------------------

\begin{definition}[Product $\sigma$-algebra]
Let $(X, \mathcal{A})$ and $(Y, \mathcal{B})$ be measurable spaces. The
\emph{product $\sigma$-algebra}\index{product $\sigma$-algebra}
$\mathcal{A} \otimes \mathcal{B}$ is the $\sigma$-algebra on $X \times Y$
generated by the \emph{measurable rectangles}:
\[
    \mathcal{A} \otimes \mathcal{B}
    = \sigma\bigl(\{A \times B : A \in \mathcal{A},\, B \in \mathcal{B}\}\bigr).
\]
\end{definition}

\begin{definition}[Sections]
For $E \subset X \times Y$ and $x \in X$, $y \in Y$, the \emph{sections}\index{section}
are:
\[
    E_x = \{y \in Y : (x,y) \in E\}, \qquad
    E^y = \{x \in X : (x,y) \in E\}.
\]
For a function $f : X \times Y \to \overline{\R}$, the sections are:
$f_x(y) = f(x,y)$ and $f^y(x) = f(x,y)$.
\end{definition}

\begin{proposition}[Measurability of sections]
\label{prop:sections_measurable}
Let $E \in \mathcal{A} \otimes \mathcal{B}$. Then for every $x \in X$, the section
$E_x \in \mathcal{B}$, and for every $y \in Y$, the section $E^y \in \mathcal{A}$.

If $f : (X \times Y, \mathcal{A} \otimes \mathcal{B}) \to
(\overline{\R}, \mathcal{B}(\overline{\R}))$ is measurable, then $f_x$ is
$\mathcal{B}$-measurable and $f^y$ is $\mathcal{A}$-measurable.
\end{proposition}

\begin{proof}
The collection $\mathcal{C} = \{E \in \mathcal{A} \otimes \mathcal{B} :
E_x \in \mathcal{B} \text{ for all } x\}$ is a $\sigma$-algebra containing the
measurable rectangles (since $(A \times B)_x = B$ if $x \in A$ and $\varnothing$
otherwise), hence $\mathcal{C} = \mathcal{A} \otimes \mathcal{B}$.
\end{proof}

% ------------------------------------------------------------
\section{Product measure}
% ------------------------------------------------------------

\begin{theorem}[Existence and uniqueness of the product measure]
\label{thm:product_measure}
\index{product measure}
Let $(X, \mathcal{A}, \mu)$ and $(Y, \mathcal{B}, \nu)$ be $\sigma$-finite measure
spaces. There exists a unique measure $\mu \otimes \nu$ on
$(X \times Y, \mathcal{A} \otimes \mathcal{B})$ such that:
\[
    (\mu \otimes \nu)(A \times B) = \mu(A) \cdot \nu(B),
    \qquad \forall A \in \mathcal{A},\, \forall B \in \mathcal{B}.
\]
Moreover, for every $E \in \mathcal{A} \otimes \mathcal{B}$:
\[
    (\mu \otimes \nu)(E) = \int_X \nu(E_x)\, d\mu(x)
    = \int_Y \mu(E^y)\, d\nu(y).
\]
\end{theorem}

\begin{proof}[Sketch of proof]
Define $\rho(E) = \int_X \nu(E_x)\, d\mu(x)$. The measurability of
$x \mapsto \nu(E_x)$ is verified by a monotone class argument. The
$\sigma$-additivity of $\rho$ follows from the monotone convergence theorem.
On rectangles, $\rho(A \times B) = \mu(A)\nu(B)$. Uniqueness follows from the
uniqueness of measures on a $\pi$-system (measurable rectangles form a
$\pi$-system) and $\sigma$-finiteness.
\end{proof}

\begin{attention}[title=\textbf{$\sigma$-finiteness hypothesis}]
$\sigma$-finiteness is essential. Without it, the product measure may not be unique,
and the Fubini and Tonelli theorems may fail.
\end{attention}

% ------------------------------------------------------------
\section{Tonelli's theorem}
% ------------------------------------------------------------

\begin{theorem}[Tonelli]
\label{thm:tonelli}
\index{Tonelli's theorem}
Let $(X, \mathcal{A}, \mu)$ and $(Y, \mathcal{B}, \nu)$ be $\sigma$-finite measure
spaces and let $f : X \times Y \to [0, +\infty]$ be
$(\mathcal{A} \otimes \mathcal{B})$-measurable. Then:
\begin{enumerate}
    \item[(i)] for $\mu$-a.e.\ $x$, the function $y \mapsto f(x,y)$ is
    $\mathcal{B}$-measurable and $x \mapsto \int_Y f(x,y)\, d\nu(y)$ is
    $\mathcal{A}$-measurable;
    \item[(ii)] the same holds with the roles of $x$ and $y$ swapped;
    \item[(iii)] we have:
    \[
        \int_{X \times Y} f\, d(\mu \otimes \nu)
        = \int_X \!\left(\int_Y f(x,y)\, d\nu(y)\right) d\mu(x)
        = \int_Y \!\left(\int_X f(x,y)\, d\mu(x)\right) d\nu(y).
    \]
\end{enumerate}
\end{theorem}

\begin{proof}
For $f = \ind_E$ with $E \in \mathcal{A} \otimes \mathcal{B}$, the result is
Theorem~\ref{thm:product_measure}. By linearity, it extends to nonnegative simple
functions. The general case follows from the monotone convergence theorem applied to
an increasing sequence of simple functions converging to $f$.
\end{proof}

% ------------------------------------------------------------
\section{Fubini's theorem}
% ------------------------------------------------------------

\begin{theorem}[Fubini]
\label{thm:fubini}
\index{Fubini's theorem}
Let $(X, \mathcal{A}, \mu)$ and $(Y, \mathcal{B}, \nu)$ be $\sigma$-finite measure
spaces and let $f \in L^1(X \times Y, \mu \otimes \nu)$. Then:
\begin{enumerate}
    \item[(i)] for $\mu$-a.e.\ $x$, $f_x \in L^1(Y, \nu)$;
    \item[(ii)] the function $x \mapsto \int_Y f(x,y)\, d\nu(y)$ (defined a.e.)
    is in $L^1(X, \mu)$;
    \item[(iii)] we have:
    \[
        \int_{X \times Y} f\, d(\mu \otimes \nu)
        = \int_X \!\left(\int_Y f(x,y)\, d\nu(y)\right) d\mu(x)
        = \int_Y \!\left(\int_X f(x,y)\, d\mu(x)\right) d\nu(y).
    \]
\end{enumerate}
\end{theorem}

\begin{proof}
Write $f = f^+ - f^-$. By Tonelli, $\int f^+\, d(\mu \otimes \nu)$ and
$\int f^-\, d(\mu \otimes \nu)$ are both finite (since $f \in L^1$). Apply Tonelli
to $f^+$ and $f^-$ separately to obtain the iterated integrals for each. In
particular, $\int_Y f^+(x,y)\, d\nu(y) < \infty$ and
$\int_Y f^-(x,y)\, d\nu(y) < \infty$ for $\mu$-a.e.\ $x$, so
$f_x \in L^1(\nu)$ a.e. Conclude by taking the difference.
\end{proof}

\begin{keyformulas}[title=\textbf{Fubini-Tonelli: user guide}]
\begin{enumerate}
    \item \textbf{Tonelli:} $f \geq 0$ measurable $\implies$ one may freely swap
    $\int\int$ (even if $= +\infty$).
    \item \textbf{Fubini:} $f \in L^1(\mu \otimes \nu)$ $\implies$ one may swap.
    \item \textbf{Strategy:} to check $f \in L^1$, first apply Tonelli to $\abs{f}$.
\end{enumerate}
\end{keyformulas}

% ------------------------------------------------------------
\section{Applications}
% ------------------------------------------------------------

\subsection{Convolution}

\begin{definition}[Convolution]
\index{convolution}
Let $f, g \in L^1(\R^n, \lambda)$ where $\lambda$ is Lebesgue measure. The
\emph{convolution} of $f$ and $g$ is:
\[
    (f * g)(x) = \int_{\R^n} f(x - y)\, g(y)\, dy.
\]
\end{definition}

\begin{proposition}
If $f, g \in L^1(\R^n)$, then $f * g$ is defined a.e., $f * g \in L^1(\R^n)$, and
$\norm{f * g}_{L^1} \leq \norm{f}_{L^1} \norm{g}_{L^1}$.
\end{proposition}

\begin{proof}
By Tonelli, $\int_{\R^n}\!\int_{\R^n} \abs{f(x-y)} \abs{g(y)}\, dy\, dx
= \norm{f}_{L^1} \norm{g}_{L^1} < \infty$ (by translation invariance).
By Fubini, the inner integral is finite for a.e.\ $x$, and the result follows.
\end{proof}

\subsection{Change of variable formula}

\begin{theorem}[Change of variables]
\index{change of variables}
Let $\varphi : U \to V$ be a $C^1$-diffeomorphism between open subsets of $\R^n$.
For every measurable $f : V \to [0, +\infty]$ (or $f \in L^1(V)$):
\[
    \int_V f(y)\, dy = \int_U f(\varphi(x))\, \abs{\det D\varphi(x)}\, dx.
\]
\end{theorem}

\begin{example}[Polar coordinates]
For $\varphi(r, \theta) = (r\cos\theta, r\sin\theta)$, we have
$\abs{\det D\varphi} = r$, giving:
\[
    \iint_{\R^2} f(x,y)\, dx\, dy
    = \int_0^{\infty}\!\int_0^{2\pi} f(r\cos\theta, r\sin\theta)\, r\, d\theta\, dr.
\]
\end{example}

\begin{example}[Gaussian integral]
Using Tonelli and polar coordinates:
\[
    \left(\int_{-\infty}^{\infty} e^{-x^2}\, dx\right)^2
    = \int_{\R^2} e^{-(x^2+y^2)}\, dx\, dy
    = \int_0^{\infty} 2\pi r\, e^{-r^2}\, dr = \pi.
\]
Hence $\int_{-\infty}^{\infty} e^{-x^2}\, dx = \sqrt{\pi}$.
\end{example}

% ------------------------------------------------------------
\section{Complements: infinite products}
% ------------------------------------------------------------

\begin{theorem}[Countable product of probability measures]
Let $(X_n, \mathcal{A}_n, \mu_n)_{n \geq 1}$ be a sequence of probability spaces.
There exists a unique probability $\PP$ on
$\bigl(\prod_{n \geq 1} X_n,\, \bigotimes_{n \geq 1} \mathcal{A}_n\bigr)$ such that:
\[
    \PP\!\left(\prod_{n=1}^{N} A_n \times \prod_{n > N} X_n\right)
    = \prod_{n=1}^{N} \mu_n(A_n)
\]
for all $N$ and all choices of $A_n \in \mathcal{A}_n$. This is the
Ionescu-Tulcea theorem (or Kolmogorov extension theorem in the product case).
\end{theorem}

% ------------------------------------------------------------
\section{Exercises}
% ------------------------------------------------------------

\begin{exercise}
Compute $\int_0^1 \int_0^1 \frac{x^2 - y^2}{(x^2 + y^2)^2}\, dx\, dy$ and
$\int_0^1 \int_0^1 \frac{x^2 - y^2}{(x^2 + y^2)^2}\, dy\, dx$. Comment on
the result. \emph{Hint}: $f \notin L^1$.
\end{exercise}

\begin{exercise}
Let $f : [0,1] \times [0,1] \to \R$ be defined by $f(x,y) = \frac{xy}{(x^2+y^2)^2}$.
Verify that $f \notin L^1([0,1]^2)$ and that the iterated integrals give different
results.
\end{exercise}

\begin{exercise}
Using Fubini, show that for $f \in L^1([0,\infty), \lambda)$:
\[
    \int_0^{\infty} \frac{1}{x} \int_0^x f(t)\, dt\, dx
    = \int_0^{\infty} f(t) \int_t^{\infty} \frac{1}{x}\, dx\, dt
\]
and explain why this equality is formal (both sides may diverge).
\end{exercise}

\begin{exercise}
Show that $\int_0^{\infty} \frac{e^{-ax} - e^{-bx}}{x}\, dx = \ln(b/a)$ for
$0 < a < b$, using Fubini.
\emph{Hint}: write $\frac{e^{-ax} - e^{-bx}}{x} = \int_a^b e^{-tx}\, dt$.
\end{exercise}

\begin{exercise}[Convolution and probability]
Let $X$ and $Y$ be independent real random variables with densities $f_X$ and $f_Y$
respectively. Show that $Z = X + Y$ has density $f_Z = f_X * f_Y$. Apply to the case
$X, Y \sim \mathcal{N}(0,1)$.
\end{exercise}

\begin{exercise}
Using polar coordinates in dimension $n$, compute the volume of the unit ball
$B_n = \{x \in \R^n : \norm{x} \leq 1\}$ and show that:
\[
    \lambda_n(B_n) = \frac{\pi^{n/2}}{\Gamma(n/2 + 1)}.
\]
\end{exercise}
